\documentclass[11pt]{article}
\usepackage[utf8]{inputenc}
\usepackage[T1]{fontenc}
\usepackage{booktabs}
\usepackage[margin=2.5cm]{geometry}
\usepackage{hyperref}

\title{Supplementary material to \textit{Measured cost of keeping one coordinate representation
per variant when ingesting a drug-resistance catalogue of the \textit{Mycobacterium
tuberculosis} complex}}
\author{Christophe Guyeux}
\date{\today}

\begin{document}
\maketitle

\section*{Supplementary Note S1: protocol of the search for a published measurement}

\textbf{Database.} A local mirror of the PubMed tuberculosis corpus, queried in read-only SQL
on 21 November 2026. Its state was measured on that day rather than taken from an earlier
record: 332\,784 records, the earliest published in April 1848 and the most recent on 20
November 2026, plus 6\,968 abstracts recovered from an open bibliographic source for records
that PubMed carries without one. Every query reads the title and the abstract in lower case,
the abstract being taken from the PubMed field or, when it is empty, from the recovered one.

\textbf{Criterion of inclusion.} A record is retained if it measures, in points of sensitivity
or specificity of a genotypic prediction of resistance, the effect of a choice of variant
representation or of coordinate conversion of a catalogue. A record that describes the
phenomenon, proposes a remedy, compares the performance of two tools without varying the
representation, or measures an effect on another criterion, is set aside with its reason.

\textbf{Queries and counts.} Seven queries were run as substring matches on the title and the
abstract. Q1, multi-nucleotide variants: 0 records. Q2, variant normalisation, representation
or left alignment: 0. Q3, HGVS notation: 1. Q4, \textit{gyrA} together with haplotype, phasing
or Ser95Thr: 14. Q5, catalogue together with coordinate or ingestion: 3. Q6, comparison or
concordance of prediction tools, naming TB-Profiler, Mykrobe, gnomonicus, PhyResSE or SAM-TB:
51. Q7, representation or notation together with sensitivity and drug resistance: 6. The 75
records were screened on their titles; none satisfies the criterion. The verbatim text of the
seven queries is given with the measurement outputs.

\textbf{Candidates read in full and reason for setting each aside.} Verboven et
al. \cite{verboven2022}, the single record of Q3: a design choice, emitting every
representation of an indel, without a measurement in points of sensitivity, and bearing on
indels where 92.2\,\% of the representations dropped here are blocks. Westhead et
al. \cite{westhead2025}: compares two implementations at a fixed representation. Hall et
al. \cite{drprg2023}: a reference-graph strategy that dissolves the question instead of
measuring it. Three further comparisons of tools (doi 10.1186/s13073-025-01458-0, five tools
on 36\,385 isolates; doi 10.1007/s00284-025-04530-5; doi 10.1371/journal.pgph.0004465, a
network meta-analysis of nine tools): compare tools, not representations. Hong et
al. \cite{hong2026}: partitions genotype-phenotype discordance without bearing on
representation, and supports the closing sentence of the Discussion instead.

\textbf{Limits of this search.} The queries match the title and the abstract, so a work that
named the phenomenon only in its body text escapes this mirror; the full-text search indexed
by Europe PMC run during an earlier pass covered that gap, but its queries were not kept
verbatim and are not reproduced here. The mirror is a tuberculosis corpus, so the works of
method published outside that field which the Discussion cites
\cite{srinivasan2021,wang2020,repun2024,tan2015} come from elsewhere. The axis of Q6, which
is the most productive of the seven, was absent from the earlier pass.

\section*{Supplementary Note S2: commits and files of the three implementations read}

The matching strategies described in the Results section on published implementations were
read in the public source code of the tools, in the state of their repositories on 16
November 2026: the library \texttt{pathogen-profiler}, which carries the matching and
annotation code of TB-Profiler, at commit \texttt{ec7dc55} (files \texttt{mutation\_db.py},
\texttt{vcf.py}, \texttt{profiler.py} and \texttt{combine\_vcf\_variants.py}); Mykrobe
\cite{hunt2019} at commit \texttt{110b4c76} (file \texttt{models.py} of the gene annotation
module); and piezo \cite{piezo} at commit \texttt{0ffab6de} (file
\texttt{grammar\_GARC1.py} and the nomenclature specification of the same repository), the
grammar being applied to genomes by gnomonicus \cite{gnomonicus}.

The four conditions that bound the codon-grouping step of TB-Profiler, each read at that
commit in the file \texttt{combine\_vcf\_}\allowbreak\texttt{variants.py}, are the following.
The step groups only
calls whose reference and alternative alleles are one base long and that fall inside a coding
exon, so nothing is grouped outside coding sequences. It counts haplotypes in the alignment
only when a BAM file carrying the MD tag is supplied, and otherwise synthesises the haplotypes
from the depth fields of the variant file. The caller named \texttt{freebayes-haplotype} sets
the alignment file to none, which removes the counting. And the whole step is skipped when the
option that splits insertions and deletions is turned off.

The probe construction of Mykrobe, read in \texttt{models.py} of its gene annotation module at
that commit, proceeds as follows. The gene model builds the inverse of the bacterial codon
table, from amino acid to the list of its codons. The construction of a coding mutation takes
the codons of the reference amino acid, uses them to remove from the alternatives the codons
shared with the reference, and then names one probe for the reference codon of the genome
against each remaining codon of the target amino acid.

\section*{Supplementary Note S3: variant calling of the genotype pipeline}

The two generations of the genotype pipeline that keep the blocks emitted by the caller
\cite{senelle2023} process reads as follows. Reads are trimmed with fastp 0.20.1, subsampled
to 200-fold coverage of the 4\,411\,532-base reference with seqkit 2.2.0, and aligned to H37Rv
(NC\_000962.3) with BWA-MEM 0.7.17 using \texttt{-Y -M}. Variants are called by snippy 4.6.0,
which clips with samclip (\texttt{-{}-max 10}) and runs freebayes with \texttt{-p 2 -P 0 -C 2
-F 0.05}, \texttt{-{}-min-coverage 10}, \texttt{-{}-min-repeat-entropy 1.0}, \texttt{-q 13}
and \texttt{-m 60}; a call is kept when its genotype is homozygous for the alternative allele,
its quality is at least 100 and its depth at least 10. The version of freebayes is the one
resolved for the batch by the package environment rather than pinned by the workflow: 1.3.6
and 1.3.2 in the batches whose logs we could audit, with samtools and bcftools 1.17 and
1.15.1 respectively. It is this caller that assembles adjacent substitutions into a single
multi-nucleotide call, and the five-base blocks of \textit{gyrA} codon 94 described in the
Results are such calls.

\section*{Supplementary Table S1}

\begin{table}[h]
\centering
\small
\caption*{\textbf{Supplementary Table S1.} The nineteen pairs of a coordinate string and a
medicine that receive both a verdict of association with resistance and a verdict of absence
of association under the catalogue that emits every representation (C1). Eleven of the
nineteen pairs concern rifampicin, and all eleven strings fall within the \textit{rpoB}
resistance-determining region; the eight others concern the two fluoroquinolones (five
pairs), ethambutol, pyrazinamide and amikacin. Seventeen distinct strings are involved, two
of them for two medicines each. Strings are in the zero-based coordinate convention of the
project, as \texttt{NC\_000962.3:position:reference:alternative}. Source:
\texttt{p77\_conflits\_C1.tsv}; the same table is provided as
\texttt{TableS1\_contradictory\_pairs\_C1.csv}.}
% Fragment généré par analyses/p712_supplementaire.py : ne pas éditer à la main.
\begin{tabular}{ll}
\toprule
Coordinate string & Medicine \\
\midrule
\texttt{NC\_000962.3:761122:GCTGTCGGGGT:CCTCGCCGGCC} & Rifampicin \\
\texttt{NC\_000962.3:761125:GTCGGGGT:CGCCGGCC} & Rifampicin \\
\texttt{NC\_000962.3:761125:GTCGGGGT:CGCCGGTC} & Rifampicin \\
\texttt{NC\_000962.3:761126:TCGGGG:GCCGGC} & Rifampicin \\
\texttt{NC\_000962.3:761126:TCGGGG:GCTGGT} & Rifampicin \\
\texttt{NC\_000962.3:761126:TCGGGGT:GCAGGTC} & Rifampicin \\
\texttt{NC\_000962.3:761126:TCGGGGT:GCCGGCC} & Rifampicin \\
\texttt{NC\_000962.3:761126:TCGGGGT:GCCGGTC} & Rifampicin \\
\texttt{NC\_000962.3:761126:TCGGGGTTG:GCTGGACTT} & Rifampicin \\
\texttt{NC\_000962.3:761137:CC:TT} & Rifampicin \\
\texttt{NC\_000962.3:761154:CGGCGCTGGGG:TGGGCCTTGGT} & Rifampicin \\
\texttt{NC\_000962.3:4247430:GGCCCGAGTC:CGCGCGGGTG} & Ethambutol \\
\texttt{NC\_000962.3:2289015:TA:GG} & Pyrazinamide \\
\texttt{NC\_000962.3:7581:ACAG:CCAC} & Moxifloxacin \\
\texttt{NC\_000962.3:7581:ACAG:GCAC} & Moxifloxacin \\
\texttt{NC\_000962.3:7567:CGC:TGT} & Levofloxacin \\
\texttt{NC\_000962.3:7581:ACAG:CCAC} & Levofloxacin \\
\texttt{NC\_000962.3:7581:ACAG:GCAC} & Levofloxacin \\
\texttt{NC\_000962.3:2715343:GTG:ATA} & Amikacin \\
\bottomrule
\end{tabular}

\end{table}

\section*{Supplementary Table S2}

\begin{table}[h]
\centering
\small
\caption*{\textbf{Supplementary Table S2.} Contribution of the decomposition of observed
blocks, in points of sensitivity, for each medicine, under the ingested catalogue C0 and under
the catalogue C1 that emits every representation, on the 50\,782 strains whose calls keep
multi-nucleotide blocks, with a 95\,\% interval from 2\,000 resamplings of sub-lineage
clusters. Outside the fluoroquinolones the contribution of the decomposition falls from 0.99
to 0.00 point for isoniazid, from 0.26 to 0.00 for rifampicin, from 0.15 to 0.00 for
streptomycin and from 0.03 to 0.00 for ethionamide. Under the ingested catalogue the blocks
that carry a determinant only through decomposition form 49 pairs of a matrix column and a
medicine, over 38 distinct columns and nine medicines. Sources: \texttt{p77\_C0.tsv} and
\texttt{p77\_C1.tsv}; the same table is provided as
\texttt{TableS2\_block\_decomposition\_gain\_by\_medicine.csv}.}
% Fragment généré par analyses/p712_supplementaire.py : ne pas éditer à la main.
\begin{tabular}{lrrrrrr}
\toprule
 & \multicolumn{3}{c}{Ingested catalogue C0} & \multicolumn{3}{c}{Catalogue C1} \\
Medicine & Gain & 95\,\% low & 95\,\% high & Gain & 95\,\% low & 95\,\% high \\
\midrule
Isoniazid & 0.99 & 0.44 & 2.32 & 0.00 & 0.00 & 0.00 \\
Rifampicin & 0.26 & 0.04 & 0.92 & 0.00 & 0.00 & 0.00 \\
Ethambutol & 0.00 & 0.00 & 0.00 & 0.00 & 0.00 & 0.00 \\
Pyrazinamide & 0.00 & 0.00 & 0.00 & 0.00 & 0.00 & 0.00 \\
Moxifloxacin & 57.37 & 29.32 & 63.33 & 13.20 & 7.86 & 15.65 \\
Levofloxacin & 49.96 & 24.08 & 57.44 & 11.75 & 6.25 & 13.31 \\
Streptomycin & 0.15 & 0.00 & 0.22 & 0.00 & 0.00 & 0.00 \\
Ethionamide & 0.03 & 0.00 & 0.16 & 0.00 & 0.00 & 0.00 \\
Amikacin & 0.00 & 0.00 & 0.00 & 0.00 & 0.00 & 0.00 \\
Kanamycin & 0.00 & 0.00 & 0.00 & 0.00 & 0.00 & 0.00 \\
Capreomycin & 0.00 & 0.00 & 0.00 & 0.00 & 0.00 & 0.00 \\
\bottomrule
\end{tabular}

\end{table}

\section*{Supplementary Table S3}

\begin{table}[h]
\centering
\small
\caption*{\textbf{Supplementary Table S3.} Sensitivity and specificity of the deterministic
baseline per medicine, with their Clopper--Pearson 95\,\% intervals, under the ingested
catalogue C0 and under the catalogue C1 that emits every representation, on the 50\,782
strains whose calls keep multi-nucleotide blocks. The lower block gives the same quantities
for the two fluoroquinolones when an observed block is matched exactly only, which are the
first two columns of Table~1 of the main text. The precision of the eleven sensitivities is
not comparable, the number of resistant strains ranging from 991 for capreomycin, whose
interval spans 0.662 to 0.721, to 13\,256 for isoniazid; rifampicin, with 10\,019 resistant
strains, has the narrowest interval, 0.926 to 0.936. Sources: \texttt{p77\_C0.tsv} and
\texttt{p77\_C1.tsv}; the same table is provided as
\texttt{TableS3\_clopper\_pearson\_intervals.csv}.}
% Fragment généré par analyses/p716_supplementaire.py : ne pas éditer à la main.
\begin{tabular}{llrrcc}
\toprule
Medicine & Ingestion & $n_R$ & $n_S$ & Sensitivity (95\,\% CI) & Specificity (95\,\% CI) \\
\midrule
Isoniazid & C0 & 13256 & 25362 & 0.903 (0.898 to 0.908) & 0.981 (0.979 to 0.983) \\
Isoniazid & C1 & 13256 & 25362 & 0.904 (0.899 to 0.909) & 0.981 (0.979 to 0.982) \\
Rifampicin & C0 & 10019 & 28665 & 0.931 (0.926 to 0.936) & 0.976 (0.974 to 0.977) \\
Rifampicin & C1 & 10019 & 28665 & 0.931 (0.926 to 0.936) & 0.976 (0.974 to 0.977) \\
Ethambutol & C0 & 5610 & 31142 & 0.830 (0.820 to 0.840) & 0.931 (0.928 to 0.933) \\
Ethambutol & C1 & 5610 & 31142 & 0.830 (0.820 to 0.840) & 0.931 (0.928 to 0.933) \\
Pyrazinamide & C0 & 2893 & 12962 & 0.733 (0.716 to 0.749) & 0.967 (0.964 to 0.970) \\
Pyrazinamide & C1 & 2893 & 12962 & 0.739 (0.723 to 0.755) & 0.967 (0.963 to 0.970) \\
Moxifloxacin & C0 & 2916 & 14931 & 0.841 (0.827 to 0.854) & 0.933 (0.929 to 0.937) \\
Moxifloxacin & C1 & 2916 & 14931 & 0.841 (0.827 to 0.854) & 0.933 (0.929 to 0.937) \\
Levofloxacin & C0 & 2358 & 10982 & 0.799 (0.783 to 0.815) & 0.971 (0.968 to 0.974) \\
Levofloxacin & C1 & 2358 & 10982 & 0.799 (0.783 to 0.815) & 0.971 (0.968 to 0.974) \\
Streptomycin & C0 & 3990 & 7642 & 0.764 (0.751 to 0.777) & 0.944 (0.939 to 0.949) \\
Streptomycin & C1 & 3990 & 7642 & 0.764 (0.751 to 0.777) & 0.944 (0.938 to 0.949) \\
Ethionamide & C0 & 3262 & 13710 & 0.711 (0.695 to 0.727) & 0.912 (0.907 to 0.917) \\
Ethionamide & C1 & 3262 & 13710 & 0.714 (0.698 to 0.730) & 0.911 (0.906 to 0.916) \\
Amikacin & C0 & 1651 & 17795 & 0.738 (0.716 to 0.759) & 0.986 (0.984 to 0.987) \\
Amikacin & C1 & 1651 & 17795 & 0.738 (0.716 to 0.759) & 0.986 (0.984 to 0.987) \\
Kanamycin & C0 & 1946 & 16998 & 0.742 (0.722 to 0.761) & 0.975 (0.973 to 0.978) \\
Kanamycin & C1 & 1946 & 16998 & 0.742 (0.722 to 0.761) & 0.975 (0.973 to 0.978) \\
Capreomycin & C0 & 991 & 8806 & 0.692 (0.662 to 0.721) & 0.976 (0.973 to 0.979) \\
Capreomycin & C1 & 991 & 8806 & 0.694 (0.664 to 0.723) & 0.976 (0.973 to 0.979) \\
\midrule
\multicolumn{6}{l}{\textit{Exact matching only (no decomposition of observed blocks)}} \\
\midrule
Moxifloxacin & C0 & 2916 & 14931 & 0.267 (0.251 to 0.284) & 0.965 (0.962 to 0.968) \\
Moxifloxacin & C1 & 2916 & 14931 & 0.709 (0.692 to 0.725) & 0.939 (0.935 to 0.943) \\
Levofloxacin & C0 & 2358 & 10982 & 0.300 (0.281 to 0.319) & 0.988 (0.985 to 0.990) \\
Levofloxacin & C1 & 2358 & 10982 & 0.682 (0.663 to 0.701) & 0.974 (0.971 to 0.977) \\
\bottomrule
\end{tabular}

\end{table}

\bibliographystyle{plain}
\bibliography{references}

\end{document}
